\begin{table}[h]
    \centering \small
        \caption{MMLU assesses a model's factual and conceptual understanding across 57 tasks spanning STEM (science, technology, engineering, mathematics), humanities, social sciences, and professional fields (e.g., law, medicine). The benchmark is commonly used to evaluate a model's ability to perform general knowledge reasoning and multitask proficiency across a diverse range of subjects and tasks. Here is an example of MMLU. }
    \label{tab:mmlu_eval_format_example}
    \begin{tabular}{p{0.95\textwidth}}
    \hline
    \textbf{PROMPT}\\
    Answer the following multiple choice question. The last line of your response should be of the following format: 'Answer: \$LETTER' (without quotes) where LETTER is one of ABCD. Think step by step before answering.\\[1ex] 
    Which tool technology is associated with Neandertals?\\[1ex]
    A. Aurignacian\\
    B. Acheulean\\
    C. Mousterian\\
    D. both b and c \\ \hline
      \textbf{Evaluation}\\
      Parse the last line in response to judge if the  choice equals to ground truth. \\ \hline
    \end{tabular}
\end{table}

\begin{table}[h]
    \centering \small
        \caption{MMLU-Redux is a subset of 5,700 manually re-annotated questions across all 57 MMLU subjects. MMLU-Redux focuses on improving the quality, clarity, and robustness of the benchmark by reducing noise, ambiguities, and potential biases in the MMLU, while potentially adjusting the scope or difficulty of tasks to better align with modern evaluation needs. Here is an example of MMLU-Redux.}
    \label{tab:mmluredux_eval_format_example}
    \begin{tabular}{p{0.95\textwidth}}
    \hline
    \textbf{PROMPT}\\
    \#\# Question: 

Sauna use, sometimes referred to as "sauna bathing," is characterized by short-term passive exposure to extreme heat \ldots In fact, sauna use has been proposed as an alternative to exercise for people who are unable to engage in physical activity due to chronic disease or physical limitations.[13]

According to the article, which of the following is NOT a benefit of sauna use?

\#\# Choices:

- (A) Decreased risk of heart attacks. \\
- (B) Increase in stroke volume. \\
- (C) Improved mental health. \\
- (D) Decreased rate of erectile dysfunction.


\#\# Instruction 

Please answer this question by first reasoning and then selecting the correct choice. \\
Present your reasoning and solution in the following json format. \\
Please show your choice in the `answer` field with only the choice letter, e.g.,`"answer": "C"`. \\

\{ \\
    "reasoning": "\_\_\_", \\
    "answer": "\_\_\_" \\
\}\\ \hline
      \textbf{Evaluation}\\
      Parse the json output in response to judge if the answer equals to ground truth. \\ \hline
    \end{tabular}
\end{table}


\begin{table}[h]
    \centering \small
        \caption{LiveCodeBench aims to evaluate model performance on the algorithm competition task, which collects new problems over time from contests across three competition platforms, namely LeetCode, AtCoder, and CodeForces. }
    \label{tab:lcb_eval_format_example}
    \begin{tabular}{p{0.95\textwidth}}
    \hline
    \textbf{PROMPT}\\
Question:
There is a stack of N cards, and the i\-th card from the top has an integer $A_i$ written on it. \\
You take K cards from the bottom of the stack and place them on top of the stack, maintaining their order. \\
Print the integers written on the cards from top to bottom after the operation. \\
Input\\
The input is given from Standard Input in the following format:\\
N K\\
$A_1 A_2 \ldots A_N$\\
Output\\
Let $B_i$ be the integer written on the i\-th card from the top of the stack after the operation. Print $B_1,B_2,\ldots,B_N$ in this order, separated by spaces.\\
Constraints\\
$-1 \leq K < N \leq 100$\\
$-1 \leq A_i \leq 100$\\
All input values are integers.\\
Sample Input 1\\
5 3 \\
1 2 3 4 5 \\
Sample Output 1 \\
3 4 5 1 2\\
Initially, the integers written on the cards are 1,2,3,4,5 from top to bottom.
After taking three cards from the bottom of the stack and placing them on top, the integers written on the cards become 3,4,5,1,2 from top to bottom.\\
\\
Sample Input 2\\
\\
6 2\\
1 2 1 2 1 2\\
\\
Sample Output 2\\
\\
1 2 1 2 1 2\\
\\
The integers written on the cards are not necessarily distinct.\\
\\
Please write a python code to solve the above problem. Your code must read the inputs from stdin and output the results to stdout.\\ \hline
      \textbf{Evaluation}\\
      Extract the code wrapped by \texttt{```}python\texttt{```} in response to judge if the answer passes the test cases. \\ \hline
    \end{tabular}
\end{table}

\begin{table}[h]
    \centering \small
        \caption{Compared to MMLU, MMLU-Pro features a curated subset of tasks, but with significantly increased difficulty. Questions in MMLU-Pro are designed to require deeper reasoning, multi-step problem-solving, and advanced domain-specific knowledge. For example, STEM tasks may involve complex mathematical derivations or nuanced scientific concepts, while humanities tasks may demand intricate contextual analysis. }
    \label{tab:mmlu-pro}
    \begin{tabular}{p{0.95\textwidth}}
    \hline
    \textbf{PROMPT}\\
    The following are multiple choice questions (with answers) about business. Think step by step and then output the answer in the format of "The answer is (X)" at the end.

\ldots

Question: Typical advertising regulatory bodies suggest, for example that adverts must not: encourage \_\_\_, cause unnecessary \_\_\_ or \_\_\_, and must not cause \_\_\_ offence. \\
Options: A. Safe practices, Fear, Jealousy, Trivial \\
B. Unsafe practices, Distress, Joy, Trivial \\
C. Safe practices, Wants, Jealousy, Trivial \\
D. Safe practices, Distress, Fear, Trivial \\
E. Unsafe practices, Wants, Jealousy, Serious \\
F. Safe practices, Distress, Jealousy, Serious \\
G. Safe practices, Wants, Fear, Serious \\
H. Unsafe practices, Wants, Fear, Trivial \\
I. Unsafe practices, Distress, Fear, Serious \\
Answer: Let's think step by step.\\ \hline
      \textbf{Evaluation}\\
      Parse the capital letter following ``Answer: '' in response to judge if the answer equals to ground truth. \\ \hline
    \end{tabular}
\end{table}

\begin{table}[h]
    \centering \small
        \caption{DROP assesses a model's ability to understand and extract relevant information from extended textual passages. Unlike simpler question-answering benchmarks that focus on factual recall, DROP requires models to process and interpret context-rich paragraphs. }
    \label{tab:drop_eval_format_example}
    \begin{tabular}{p{0.95\textwidth}}
    \hline
    \textbf{PROMPT}\\
    You will be asked to read a passage and answer a question. Some examples of passages and Q\&A are provided below. 

\# Examples  
---
Passage:  Looking to avoid back-to-back divisional losses, the Patriots traveled to Miami to face the 6-4 Dolphins at Dolphin Stadium \ldots Cassel's 415 passing yards made him the second quarterback in Patriots history to throw for at least 400 yards in two or more games; Drew Bledsoe had four 400+ yard passing games in his Patriots career.

Question: How many points did the Dolphins lose by?
Answer: 20.

---
Passage:  In week 2, the Seahawks took on their division rivals, the San Francisco 49ers. Prior to the season, NFL analysts rated this rivalry as the top upcoming rivalry, as well as the top rivalry of the decade \ldots Seattle was now 2-0, and still unbeaten at home.

Question: How many field goals of at least 30 yards did Hauschka make?
Answer: 2.

---
Passage:  at Raymond James Stadium, Tampa, Florida TV Time: CBS 1:00pm eastern The Ravens opened the regular season on the road against the Tampa Bay Buccaneers on September 10. \ldots With the win, the Ravens were 1-0 and 1-0 against NFC Opponents.

Question: how many yards did lewis get
Answer: 4.
\# Your Task

---
Passage:  The Chargers (1-0) won their season opener 22-14 against the Oakland Raiders after five field goals by Nate Kaeding and three botched punts by the Raiders. The Raiders Pro Bowl long snapper Jon Condo suffered a head injury in the second quarter. He was replaced by linebacker Travis Goethel, who had not snapped since high school. Goethel rolled two snaps to punter Shane Lechler, each giving the Chargers the ball in Raiders territory, and Lechler had another punt blocked by Dante Rosario. The Chargers scored their only touchdown in the second quarter after a 13-play, 90-yard drive resulted in a 6-yard touchdown pass from Philip Rivers to wide receiver Malcom Floyd. The Chargers failed to score four out of five times in the red zone. San Diego led at halftime 10-6, and the Raiders did not scored a touchdown until 54 seconds remained in the game. Undrafted rookie Mike Harris made his first NFL start, filing in for left tackle for an injured Jared Gaither. San Diego protected Harris by having Rivers throw short passes; sixteen of Rivers' 24 completions were to running backs and tight ends, and he threw for 231 yards while only being sacked once. He did not have an interception after throwing 20 in 2011. The win was the Chargers' eighth in their previous nine games at Oakland. It improved Norv Turner's record to 4-2 in Chargers' season openers. Running back Ryan Mathews and receiver Vincent Brown missed the game with injuries.

Question: How many yards did Rivers pass?
Answer: 

Think step by step, then write a line of the form "Answer: \$ANSWER" at the end of your response. \\ \hline
      \textbf{Evaluation}\\
      Parse the capital letter following ``Answer: '' in response to judge if the answer equals to ground truth. \\ \hline
    \end{tabular}
\end{table}

\begin{table}[h]
    \centering \small
        \caption{Instruction‑Following Evaluation (IFEval) is a benchmark designed to assess a model’s ability to comply with explicit, verifiable instructions embedded within prompts. It targets a core competency of large language models (LLMs): producing outputs that meet multiple, clearly defined constraints specified by the user.}
    \label{tab:ifeval_eval_format_example}
    \begin{tabular}{p{0.95\textwidth}}
    \hline
    \textbf{PROMPT}\\
    Kindly summarize the text below in XML format. Make sure the summary contains less than 4 sentences.

Quantum entanglement is the phenomenon that occurs when a group of particles are generated, interact, or share spatial proximity in such a way that the quantum state of each particle of the group cannot be described independently of the state of the others, including when the particles are separated by a large distance. The topic of quantum entanglement is at the heart of the disparity between classical and quantum physics: entanglement is a primary feature of quantum mechanics not present in classical mechanics.

Measurements of physical properties such as position, momentum, spin, and polarization performed on entangled particles can, in some cases, be found to be perfectly correlated. For example, if a pair of entangled particles is generated such that their total spin is known to be zero, and one particle is found to have clockwise spin on a first axis, then the spin of the other particle, measured on the same axis, is found to be anticlockwise. However, this behavior gives rise to seemingly paradoxical effects: any measurement of a particle's properties results in an apparent and irreversible wave function collapse of that particle and changes the original quantum state. With entangled particles, such measurements affect the entangled system as a whole.

Such phenomena were the subject of a 1935 paper by Albert Einstein, Boris Podolsky, and Nathan Rosen, and several papers by Erwin Schrödinger shortly thereafter, describing what came to be known as the EPR paradox. Einstein and others considered such behavior impossible, as it violated the local realism view of causality (Einstein referring to it as "spooky action at a distance") and argued that the accepted formulation of quantum mechanics must therefore be incomplete.\\ \hline
      \textbf{Evaluation}\\
      Call official functions to check if the answer is consistent with the instructions. \\ \hline
    \end{tabular}
\end{table}

\begin{table}[h]
    \centering \small
        \caption{FRAMES (Factuality, Retrieval, And reasoning MEasurement Set) is a comprehensive benchmark designed to evaluate core components of retrieval-augmented generation (RAG) systems. Our evaluation employs the benchmark's official "Oracle Prompt" configuration. In this setting, each test prompt includes the question along with all the ground truth Wikipedia articles, thus eliminating the need for an external retrieval component (e.g., BM25). This setting allows us to specifically measure a model's ability to reason over and synthesize information from provided sources to generate correct and verifiable facts. }
    \label{tab:frames_eval_format_example}
    \begin{tabular}{p{0.95\textwidth}}
    \hline
    \textbf{PROMPT}\\
    Here are the relevant Wikipedia articles: \\
url: https:\/\/en.wikipedia.org\/wiki\/President\_of\_the\_United\_States \\
url content: The president of the United States (POTUS) is the head of state and head of government of the United States of America. The president directs the executive branch of the federal government and is the commander-in-chief of the United States Armed Forces.
\ldots

Based on all the information, answer the query. 

Query: If my future wife has the same first name as the 15th first lady of the United States' mother and her surname is the same as the second assassinated president's mother's maiden name, what is my future wife's name? \\ \hline
      \textbf{Evaluation}\\
      ===Task===\\
I need your help in evaluating an answer provided by an LLM against a ground truth answer. Your task is to determine if the ground truth answer is present in the LLM's response. Please analyze the provided data and make a decision. \\
===Instructions=== \\
1. Carefully compare the "Predicted Answer" with the "Ground Truth Answer".\\
2. Consider the substance of the answers - look for equivalent information or correct answers.\\
Do not focus on exact wording unless the exact wording is crucial to the meaning.\\
3. Your final decision should be based on whether the meaning and the vital facts of the
"Ground Truth Answer" are present in the "Predicted Answer:"\\
===Input Data===
- Question: If my future wife has the same first name as the 15th first lady of the United States' mother and her surname is the same as the second assassinated president's mother's maiden name, what is my future wife's name? \\
- Predicted Answer: \ldots \\
- Ground Truth Answer: Jane Ballou\\
===Output Format===\\
Provide your final evaluation in the following format:\\
Explanation: xxx\\
Decision: "TRUE" or "FALSE"\\
Please proceed with the evaluation.\\ \hline
    \end{tabular}
\end{table}


\begin{table}[h]
    \centering \small
        \caption{Arena-Hard is an open-ended evaluation benchmark specifically designed to assess the capabilities of LLMs. It presents models with challenging, novel, and diverse prompts curated from Chatbot Arena, a continuously evolving, crowd-sourced platform. It focuses on measuring model performance in open-ended tasks, with particular emphasis on coding and mathematics-related prompts.
        Given the inherently subjective nature of open-ended tasks, where multiple valid responses may exist, the benchmark necessitates the use of an evaluation model to approximate human judgment effectively. Higher evaluation scores suggest that the model is more likely to be favored by human users in real-world scenarios.   }
    \label{tab:arenahard_eval_format_example}
    \begin{tabular}{p{0.95\textwidth}}
    \hline
    \textbf{PROMPT}\\
    To write an SQL query that selects the top 10 rows in a database and joins to 3 different tables based on a field called "code", you would need to know the names of the tables and the specific fields you want to select. Assuming the tables are named "table1", "table2", and "table3", and you want to select all fields from each table, the query would look like this:

```sql
SELECT *
FROM (
  SELECT * 
  FROM table1
  LIMIT 10
) AS t1
JOIN table2 ON t1.code = table2.code
JOIN table3 ON t1.code = table3.code;
```

This query first selects the top 10 rows from "table1" and then joins "table2" and "table3" based on the "code" field. If you want to select specific fields from each table, you can replace the asterisks (*) with the field names separated by commas.\\ \hline
      \textbf{Evaluation}\\
      Please act as an impartial judge and evaluate the quality of the responses provided by two AI assistants to the user prompt displayed below. You will be given assistant A's answer and assistant B's answer. Your job is to evaluate which assistant's answer is better.

Begin your evaluation by generating your own answer to the prompt. You must provide your answers before judging any answers.

When evaluating the assistants' answers, compare both assistants' answers with your answer. You must identify and correct any mistakes or inaccurate information.

Then consider if the assistant's answers are helpful, relevant, and concise. Helpful means the answer correctly responds to the prompt or follows the instructions. Note when user prompt has any ambiguity or more than one interpretation, it is more helpful and appropriate to ask for clarifications or more information from the user than providing an answer based on assumptions. Relevant means all parts of the response closely connect or are appropriate to what is being asked. Concise means the response is clear and not verbose or excessive.

Then consider the creativity and novelty of the assistant's answers when needed. Finally, identify any missing important information in the assistants' answers that would be beneficial to include when responding to the user prompt.

After providing your explanation, you must output only one of the following choices as your final verdict with a label:

1. Assistant A is significantly better: [[A$>>$B]]\\
2. Assistant A is slightly better: [[$A>$B]]\\
3. Tie, relatively the same: [[A$=$B]]\\
4. Assistant B is slightly better: [[B$>$A]]\\
5. Assistant B is significantly better: [[B$>>$A]]\\

Example output: "My final verdict is tie: [[A$=$B]]". \\ \hline
    \end{tabular}
\end{table}


\begin{table}[h]
    \centering \small
        \caption{AlpacaEval 2.0 is an open-ended evaluation dataset, similar in nature to ArenaHard, and leverages an LLM to assess model performance on subjective tasks. However, in contrast to ArenaHard, the prompts in AlpacaEval 2.0 are generally less challenging and only a small subset necessitates the deployment of reasoning capabilities by the evaluated models. }
    \label{tab:alpacaeval_eval_format_example}
    \begin{tabular}{p{0.95\textwidth}}
    \hline
    \textbf{PROMPT}\\
    What are the names of some famous actors that started their careers on Broadway?
    \\ \hline
      \textbf{Evaluation}\\
      $<|im_start|>$system \\
You are a highly efficient assistant, who evaluates and selects the best large language model (LLMs) based on the quality of their responses to a given instruction. This process will be used to create a leaderboard reflecting the most accurate and human-preferred answers.\\
$<|im_end|>$\\
$<|im_start|>$user\\
I require a leaderboard for various large language models. I'll provide you with prompts given to these models and their corresponding outputs. Your task is to assess these responses, and select the model that produces the best output from a human perspective.

\#\# Instruction

\{\\
    "instruction": """\{instruction\}""",\\
\}\\

\#\# Model Outputs

Here are the unordered outputs from the models. Each output is associated with a specific model, identified by a unique model identifier.

\{\\
    \{\\
        "model\_identifier": "m",\\
        "output": """\{output\_1\}"""\\
    \},\\
    \{\\
        "model\_identifier": "M",\\
        "output": """\{output\_2\}"""\\
    \}\\
\}

\#\# Task

Evaluate the models based on the quality and relevance of their outputs, and select the model that generated the best output. Answer by providing the model identifier of the best model. We will use your output as the name of the best model, so make sure your output only contains one of the following model identifiers and nothing else (no quotes, no spaces, no new lines, ...): m or M.

\#\# Best Model Identifier\\
$<|im_end|>$\\ \hline
    \end{tabular}
\end{table}

\begin{table}[h]
    \centering \small
        \caption{The CLUEWSC (Chinese Language Understanding Evaluation Benchmark - Winograd Schema Challenge) is a specialized task within the CLUE benchmark suite designed to evaluate a model's commonsense reasoning and contextual understanding capabilities in Chinese. }
    \label{tab:cluewsc_eval_format_example}
    \begin{tabular}{p{0.95\textwidth}}
    \hline
    \textbf{PROMPT}\\
    请参考示例的格式，完成最后的测试题。

下面是一些示例：
他伯父还有许多女弟子，大半是富商财主的外室；这些财翁白天忙着赚钱，怕小公馆里的情妇长日无聊，要不安分，常常叫她们学点玩艺儿消遣。\\
上面的句子中的"她们"指的是\\
情妇

耶律克定说到雁北义军时，提起韦大哥，就连声说不可挡、不可挡，似有谈虎色变之味。后来又听说粘罕在云中，特派人厚币卑词，要与‘韦义士修好’。韦大哥斩钉截铁地回绝了，大义凛然，端的是条好汉。如今张孝纯也想结识他，几次三番派儿子张浃上门来厮缠，定要俺引他上雁门山去见韦大哥。\\
上面的句子中的"他"指的是\\
张浃

“你何必把这事放在心上？何况你的还不过是手稿，并没有发表出来。”龙点睛越发坦率：“如果发表出来了，那倒也就算了。不过既然没发表出来，我何必还让它飘在外头呢？你给我找一找吧，我要收回。”\\
上面的句子中的"它"指的是 \\
手稿

这个身材高大，曾经被孙光平拿着菜刀追赶得到处乱窜的年轻人，那天早晨穿上了全新的卡其布中山服，像一个城里来的干部似的脸色红润，准备过河去迎接他的新娘。\\
上面的句子中的"他"指的是 \\
年轻人

负责接待我们的是两位漂亮的朝鲜女导游，身材高挑，露出比例完美的小腿。一个姓韩，一个姓金，自称小韩和小金。她们的中文说得毫无口音，言谈举止也相当亲切。\\
上面的句子中的"她们"指的是 \\
两位漂亮的朝鲜女导游

下面是测试题，请在思考结束后（</think>后）用一句话输出答案，不要额外的解释。

崩龙珍夫妻康健和美；鞠琴十年前丧偶，两年前重结良缘，现在的老伴是一位以前未曾有过婚史的高级工程师；崩龙珍和鞠琴都尽量避免谈及自己的爱人，也尽量回避提及蒋盈波的亡夫屈晋勇——尽管她们对他都很熟悉；当然也绝不会愚蠢地提出蒋盈波今后是一个人过到底还是再找个老伴的问题来加以讨论，那无论如何还为时过早。\\
上面的句子中的"他"指的是
\\ \hline
      \textbf{Evaluation}\\
      Parse the last line in response to judge if the answer equals to ground truth.
       \\ \hline
    \end{tabular}
\end{table}


\begin{table}[h]
    \centering \small
        \caption{C-EVAL evaluates a model's breadth and depth of knowledge across 52 diverse academic disciplines, spanning humanities, social sciences, STEM (Science, Technology, Engineering, and Mathematics), and other professional fields (e.g., medicine, law). All question in C-Eval are Chinese. }
    \label{tab:ceval_eval_format_example}
    \begin{tabular}{p{0.95\textwidth}}
    \hline
    \textbf{PROMPT}\\
以下是中国关于逻辑学考试的单项选择题，请选出其中的正确答案。

1991年6月15日，菲律宾吕宋岛上的皮纳图博火山突然大喷发，2000万吨二氧化硫气体冲入平流层，形成的霾像毯子一样盖在地球上空，把部分要照射到地球的阳光反射回太空几年之后，气象学家发现这层霾使得当时地球表面的温度累计下降了0．5℃，而皮纳图博火山喷发前的一个世纪，因人类活动而造成的温室效应已经使地球表面温度升高1℃。某位持“人工气候改造论”的科学家据此认为，可以用火箭弹等方式将二氧化硫充入大气层，阻挡部分阳光，达到地球表面降温的目的。以下哪项如果为真,最能对该科学家的提议构成质疑?\_\_\_ \\
A. 如果利用火箭弹将二氧化硫充入大气层，会导致航空乘客呼吸不适。\\
B. 火山喷发形成的降温效应只是暂时的，经过一段时间温度将再次回升。\\
C. 可以把大气层中的碳取出来存储在地下，减少大气层的碳含量。\\
D. 不论何种方式，“人工气候改造”都将破坏地区的大气层结构。\\
答案： B

\ldots

新疆的哈萨克人用经过训练的金雕在草原上长途追击野狼。某研究小组为研究金雕的飞行方向和判断野狼群的活动范围，将无线电传导器放置在一只金雕身上进行追踪。野狼为了觅食，其活动范围通常很广。因此，金雕追击野狼的飞行范围通常也很大。然而两周以来，无线电传导器不断传回的信号显示，金雕仅在放飞地3公里的范围内飞行。以下哪项如果为真，最有助于解释上述金雕的行为?\_\_\_ \\
A. 金雕放飞地周边重峦叠嶂，险峻异常。\\
B. 金雕的放飞地2公里范围内有一牧羊草场，成为狼群袭击的目标。\\
C. 由于受训金雕的捕杀，放飞地广阔草原的野狼几乎灭绝了。\\
D. 无线电传导信号仅能在有限的范围内传导。\\
\\ \hline
      \textbf{Evaluation}\\
      Parse the last line in response to judge if the choice equals to ground truth.
       \\ \hline
    \end{tabular}
\end{table}




\begin{table}[h]
    \centering \small
        \caption{GPQA (Graduate‑Level Google‑Proof QA Benchmark) is a rigorous evaluation framework designed to measure an LLM’s ability to tackle complex, graduate-level multiple‑choice problems in STEM domains—specifically biology, physics, and chemistry. }
    \label{tab:gpqa_eval_format_example}
    \begin{tabular}{p{0.95\textwidth}}
    \hline
    \textbf{PROMPT}\\
    Answer the following multiple choice question. The last line of your response should be of the following format: 'ANSWER: \$LETTER' (without quotes) where LETTER is one of ABCD. Think step by step before answering.

Two quantum states with energies E1 and E2 have a lifetime of $10^{-9}$ sec and $10^{-8}$ sec, respectively. We want to clearly distinguish these two energy levels. Which one of the following options could be their energy difference so that they can be clearly resolved?

A) $10^{-9}$ eV \\
B) $10^{-8}$ eV \\
C) $10^{-4}$ eV \\
D) $10^{-11}$ eV \\ \hline
      \textbf{Evaluation}\\
      Parse the capital letter following ``ANSWER: '' in response to judge if the choice equals to ground truth. \\ \hline
    \end{tabular}
\end{table}

\begin{table}[h]
    \centering \small
        \caption{SimpleQA is a factuality evaluation benchmark that measures a model’s ability to answer short, fact-seeking questions with precise, verifiable correctness. }
    \label{tab:simpleqa_eval_format_example}
    \begin{tabular}{p{0.95\textwidth}}
    \hline
    \textbf{PROMPT}\\
    Who received the IEEE Frank Rosenblatt Award in 2010? \\ \hline
      \textbf{Evaluation}\\
      Your job is to look at a question, a gold target, and a predicted answer, and then assign a grade of either ["CORRECT", "INCORRECT", "NOT\_ATTEMPTED"].
First, I will give examples of each grade, and then you will grade a new example.


The following are examples of CORRECT predicted answers.

Question: What are the names of Barack Obama's children? \\
Gold target: Malia Obama and Sasha Obama \\
Predicted answer 1: sasha and malia obama \\
Predicted answer 2: most people would say Malia and Sasha, but I'm not sure and would have to double check \\
\ldots


The following are examples of INCORRECT predicted answers.

Question: What are the names of Barack Obama's children? \\
Gold target: Malia and Sasha \\
Predicted answer 1: Malia. \\
Predicted answer 2: Malia, Sasha, and Susan. \\
\ldots


The following are examples of NOT\_ATTEMPTED predicted answers.

Question: What are the names of Barack Obama's children? \\
Gold target: Malia and Sasha \\
Predicted answer 1: I don't know. \\
Predicted answer 2: I need more context about which Obama you are 
\ldots


Also note the following things: \\
\ldots

Here is a new example. Simply reply with either CORRECT, INCORRECT, NOT\_ATTEMPTED. Don't apologize or correct yourself if there was a mistake; we are just trying to grade the answer.\\

Question: Who received the IEEE Frank Rosenblatt Award in 2010?\\
Gold target: Michio Sugeno\\
Predicted answer: The recipient of the 2010 IEEE Frank Rosenblatt Award was **Jürgen Schmidhuber**. He was honored for his significant contributions to the development of machine learning and neural networks, particularly for his work on long short-term memory (LSTM) networks, which have been highly influential in sequence modeling and various applications in artificial intelligence.


Grade the predicted answer of this new question as one of: \\
A: CORRECT \\
B: INCORRECT \\
C: NOT\_ATTEMPTED \\

Just return the letters "A", "B", or "C", with no text around it. \\ \hline
    \end{tabular}
\end{table}


\begin{table}[h]
    \centering \small
        \caption{An example of C-SimpleQA. It measures a model’s ability to answer short, fact-seeking questions in Chinese with precise, verifiable correctness. }
    \label{tab:csimpleqa_eval_format_example}
    \begin{tabular}{p{0.95\textwidth}}
    \hline
    \textbf{PROMPT}\\
显脉香茶菜可以用来治疗急性的什么类型的黄疸型肝炎？
\\ \hline
      \textbf{Evaluation}\\
      请根据给定问题、标准答案和模型预测的答案来评估模型的回答是否正确。您的任务是将结果评定为：【正确】、【错误】或【未尝试】。

首先，我们将列出每个评定类别的示例，然后请您对新问题的预测答案进行评定。

以下是【正确】的答复示例：

问题：贝拉克·奥巴马的孩子叫什么名字？\\
标准答案：玛丽亚·奥巴马和萨莎·奥巴马\\
模型预测1：Malia Obama and Sasha Obama\\
模型预测2：玛丽亚和萨沙\\
\ldots

以下是【错误】的答复示例：

问题：巴拉克·奥巴马的孩子叫什么名字？\\
标准答案：玛丽亚·奥巴马和萨莎·奥巴马\\
模型预测1：玛丽亚\\
模型预测2：玛丽亚、萨莎和苏珊\\
\dots

以下是【未尝试】的答复示例：

问题：巴拉克·奥巴马的孩子叫什么名字？\\
标准答案：玛丽亚·奥巴马和萨莎·奥巴马\\
模型预测1：我不知道。\\
模型预测2：我需要更多关于您所指奥巴马的上下文。\\
\ldots

下面是一个新的问题示例。请只回复A、B、C之一，不要道歉或纠正自己的错误，只需要评估该回答。

问题: 显脉香茶菜可以用来治疗急性的什么类型的黄疸型肝炎？\\
正确答案: 黄疸型肝炎\\
预测答案: \dots


将此新问题的预测答案评定为以下之一：\\
A:【正确】\\
B:【错误】\\
C:【未尝试】\\

只返回字母"A"、"B"或"C"，无须添加其他文本。
       \\ \hline
    \end{tabular}
\end{table}

% =============== math ===============

\begin{table}[h]
    \centering \small
        \caption{An example of math evaluation, which applies to AIME, MATH, and CNMO. These benchmarks evaluate model performance on mathematical tasks. } 
    \label{tab:math_eval_format_example}
    \begin{tabular}{p{0.95\textwidth}}
    \hline
    \textbf{PROMPT}\\
    Let $b \geq 2$ be an integer. Call a positive integer $n$ $b\textit{-eautiful}$ if it has exactly two digits when expressed in base $b$, and these two digits sum to $\sqrt{n}$. For example, $81$ is $13$-eautiful because $81=\underline{6}$ $\underline{3}_{13}$ and $6+3=\sqrt{81}$. Find the least integer $b\geq 2$ for which there are more than ten $b$-eautiful integers. \\
    Please reason step by step, and put your final answer within \textbackslash boxed\{\}. \\
    \hline
  \textbf{Evaluation}\\
  
  Parse the final answer within \textbackslash boxed\{\} and use a rule-based grader to determine if it equals the ground truth. Round numerical values as needed, and use `SymPy'\footnote{\url{https://www.sympy.org}} to parse expressions. \\
      \hline
    \end{tabular}
\end{table}